\Article{25}
Um auch den anderen Teil unserer Theorie durch ein Beispiel zu veranschaulichen, untersuchen wir die Genauigkeit, mit der die Seite Falkenberg–Breithorn mithilfe ausgeglichener Beobachtungen aus der Seite Wilsede–Wulfsode bestimmt wird. Die Funktion \(u\), durch die sie in diesem Fall ausgedrückt wird, lautet
\SourceLeaf{097}
\[
u=22877^{\mathrm m}94\times\frac{\sin(v^{(13)}-v^{(12)}-0''652).\sin(v^{(14)}-v^{(16)}-0''814)}{\sin(v^{(1)}-v^{(0)}-0''652).\sin(v^{(6)}-v^{(4)}-0''814)}
\]

Ihr Wert ergibt sich aus den verbesserten Werten der Richtungen \(v^{(0)},v^{(1)}\) usw. zu
\[
=26766^{\mathrm m}68
\]

Die Differentiation dieses Ausdrucks ergibt, wenn die Differentiale \(\mathrm dv^{(0)},\allowbreak\mathrm dv^{(1)}\) usw. in Bogensekunden ausgedrückt gedacht werden,
\[
\begin{aligned}
\mathrm du={}&0^{\mathrm m}16991(\mathrm dv^{(0)}-\mathrm dv^{(1)})+0^{\mathrm m}08836(\mathrm dv^{(4)}-\mathrm dv^{(6)})\\
&-0^{\mathrm m}03899(\mathrm dv^{(12)}-\mathrm dv^{(13)})+0^{\mathrm m}16731(\mathrm dv^{(14)}-\mathrm dv^{(16)})
\end{aligned}
\]

Hieraus ergibt sich ferner
\[
\begin{aligned}\relax
[al]&=-0{,}08836\\
[bl]&=+0{,}13092\\
[cl]&=-0{,}00260\\
[dl]&=+0{,}07895\\
[el]&=+0{,}03899\\
[fl]&=-40{,}1315\\
[gl]&=+10{,}9957\\
[ll]&=+0{,}13238
\end{aligned}
\]

Schließlich ergibt sich nach den oben dargelegten Methoden, wenn wir den Meter als Einheit der linearen Maße annehmen,
\[
\frac1P=0{,}08329,\quad\text{oder}\quad P=12{,}006
\]
woraus sich für den Wert der Seite Falkenberg–Breithorn ein zu erwartender mittlerer Fehler von \(=0{,}2886m\) Metern ergibt (wobei \(m\) der zu erwartende mittlere Fehler der beobachteten Richtungen, ausgedrückt in Bogensekunden, ist); nehmen wir daher den oben ermittelten Wert von \(m\) an,
\[
=0^{\mathrm m}1209
\]

Im Übrigen zeigt die Betrachtung des Dreieckssystems unmittelbar, dass der Punkt Hauselberg vollständig daraus hätte weggelassen werden können, ohne die Verbindung zwischen den Seiten Wilsede–Wulfsode und Falkenberg–Breithorn zu beeinträchtigen. Es widerspricht jedoch einer guten Methode, deshalb Beobachtungen zu \textit{unterdrücken}, die sich auf den Punkt Hauselberg beziehen\SourceLeaf{098}\footnote{Der größte Teil dieser Beobachtungen war bereits gemacht worden, bevor der Punkt Breithorn gefunden und in das System aufgenommen wurde.}, da sie sicherlich zur Erhöhung der Genauigkeit beitragen können. Um deutlicher zu zeigen, wie groß der dadurch erzielte Genauigkeitszuwachs ist, haben wir die Rechnung wiederholt und alles ausgeschlossen, was sich auf den Punkt Hauselberg bezieht; dabei entfallen acht der oben angegebenen 18 Richtungen, und die wahrscheinlichsten Fehler der übrigen ergeben sich wie folgt:
\[
\begin{array}{r@{\;}l@{\qquad}|@{\qquad}r@{\;}l}
(0)=&+0''327&(12)=&+0''206\\
(1)=&-0{,}206&(13)=&-0{,}206\\
(3)=&-0{,}121&(14)=&+0{,}327\\
(4)=&+0{,}121&(15)=&+0{,}206\\
(6)=&-0{,}121&(16)=&+0{,}121
\end{array}
\]
Der Wert der Seite Falkenberg–Breithorn ergibt sich dann zu \(=26766^{\mathrm m}63\) und weicht somit nur wenig von dem oben ermittelten Wert ab; die Gewichtsrechnung ergibt jedoch
\[
\frac1P=0{,}13082\quad\text{oder}\quad P=7{,}644
\]
und somit beträgt der zu erwartende mittlere Fehler \(=0{,}36169m\) Meter \(=0^{\mathrm m}1515\). Daraus ist ersichtlich, dass durch Hinzunahme der Beobachtungen; die sich auf den Punkt Hauselberg beziehen, das Gewicht der Bestimmung der Seite Falkenberg–Breithorn im Verhältnis von 7,644 zu 12,006, also von der Einheit zu 1.571, zugenommen hat.
\par\medskip\begin{center}\rule{.28\textwidth}{.3pt}\end{center}
